\documentclass[12pt]{article}

\usepackage[margin=1in]{geometry}
\usepackage{amsmath, amssymb}
\usepackage{graphicx}
\usepackage{booktabs}
\usepackage{caption}
\usepackage{subcaption}
\usepackage{float}
\usepackage{microtype}
\usepackage[hidelinks]{hyperref}
\usepackage{parskip}
\usepackage{hyperref}
\setlength{\parindent}{0pt}
\setlength{\parskip}{6pt}

\graphicspath{
  {outputs/damped_oscillator/}
  {outputs/lorenz/bic_standard/}
  {outputs/cylinder_wake/bic_standard/}
}

\captionsetup{font=small, labelfont=bf, skip=4pt}
\captionsetup[sub]{font=small, labelfont=bf}

\title{\textbf{ME 697 Final Report}\\[4pt]
\large Discovering Governing Equations from Data:\\
Sparse Identification of Nonlinear Dynamical Systems}
\author{Kylie Sommer-Kohrt}
\date{5/9/2026}

\begin{document}
\maketitle
\thispagestyle{empty}
% ============================================================
\section*{1. Literature Review}
% ============================================================

\section{The Shift Toward Physics-Informed Machine Learning}
System identification has undergone a fundamental transformation since its inception in the 1950s, evolving from purely statistical and frequency-domain methods toward the modern paradigm of Physics-Informed Machine Learning (PiML). Historically, the field was divided between rigid, physics-heavy models that struggled with unmodeled dynamics and high-dimensional, black-box models that lacked interpretability and generalizability. SINDy (Sparse Identification of Nonlinear Dynamics) emerged as a definitive bridge in this landscape, positioning itself as a method that enforces physical consistency through the mathematical lens of parsimony rather than just data-fitting \cite{Brunton2016}. As noted in the broader literature regarding PiML status and challenges, this shift represents a move toward architectures that are both data-efficient and physically grounded \cite{Zhang2025}.

\section{Historical Lineage and Precursors to SINDy}
What led to the development of SINDy was a convergence of two major mathematical lineages: Koopman operator theory and the revolution in compressed sensing. In the years preceding the seminal work by Brunton et al., the research community relied heavily on Dynamic Mode Decomposition (DMD) to extract spatial-temporal patterns from data. While DMD was highly effective for linear approximations, it fell short in capturing the full complexity of nonlinear dynamical systems. Researchers looked toward the Koopman operator, which suggests that nonlinear dynamics can be represented linearly in an infinite-dimensional space of observables \cite{Kutz2023}. SINDy solved the practical challenge of finding a finite lifting set by re-envisioning the governing equations as a sparse signal within a high-dimensional library of candidate functions. By applying sparse regression techniques, researchers could "select" the correct physical terms from a dictionary, essentially automating symbolic discovery \cite{Brunton2016}.

\section{SINDy’s Place in the PiML Paradigm}
Within the PiML framework, SINDy occupies a distinct niche compared to other methods like Physics-Informed Neural Networks (PINNs). While PINNs use the known residuals of differential equations to regularize the training of a neural network, SINDy discovers the differential equations themselves, providing a high degree of structural freedom while maintaining a strong physics prior. This interpretability is crucial for safety-critical fields such as aerospace, where identified coefficients must align with known aerodynamic stability derivatives. This methodology allows practitioners to absorb deviations from nominal physics into the learned structure rather than treating them solely as bias \cite{Kutz2023}. Furthermore, the convergence of these sparse identification algorithms has been rigorously analyzed to ensure they remain robust in the presence of noise and limited data \cite{Zhang2019}.

\section{Promising Research Developments (2024--2026)}
As of early 2026, the research frontier has moved beyond the original "fixed library" constraint. One of the most promising developments is the integration of SINDy with Kolmogorov-Arnold Networks (KANs). SINDy-KANs utilize learnable spline-based functions on edges, allowing the model to discover optimal functional forms on-the-fly without requiring a pre-defined exhaustive library \cite{Vaddireddy2026}. This effectively mitigates the "library bottleneck" previously cited as a primary drawback of the method.

Another significant advancement is the development of the Library Optimization Mechanism (SINDy-LOM), which introduces a two-layer optimization architecture to improve recursive long-term prediction accuracy \cite{Takeishi2026}. Additionally, the SINDy-SHRED (Shallow Recurrent Decoders) framework has emerged to solve the challenge of "hidden" states and sensor sparsity. By merging SINDy-class dynamics with recurrent architectures, researchers can map temporal sequences from sparse sensors to a latent space where the governing laws are symbolically identified \cite{Morales2026}. These hybrid architectures represent the state-of-the-art in 2026, providing a robust pathway for real-time, adaptive modeling in complex aerospace environments.
% ============================================================
\section*{2. Summary of Paper Contributions}
% ============================================================

Brunton et al.\ (2016) make four interconnected contributions.

\textbf{The SINDy framework.}
The central insight is to recast nonlinear system identification as sparse
regression against a user-defined library $\boldsymbol{\Theta}(\mathbf{X})$.
Prior approaches either (i) assumed linearity, (ii) searched over model
structures combinatorially, or (iii) produced dense coefficient vectors with
no physical interpretability.  By writing $\dot{\mathbf{X}} =
\boldsymbol{\Theta}(\mathbf{X})\boldsymbol{\Xi}$ and seeking a sparse
$\boldsymbol{\Xi}$, SINDy simultaneously identifies \emph{which} terms are
active and estimates their coefficients, all within a single regression
step.  The recovered equations are human-readable and, for physical systems
whose true governing laws are polynomials or trigonometric combinations
(fluid mechanics, structural dynamics, chemical kinetics), can match the
ground truth exactly.

\textbf{Sequential Thresholded Least Squares (STLSQ).}
The authors introduce STLSQ as a practical sparse solver for the SINDy
regression.  Unlike LASSO, STLSQ enforces sparsity via a hard threshold
rather than an $\ell_1$ penalty, which eliminates the bias that $\ell_1$
regularisation imposes on retained coefficients.  The algorithm alternates
between hard-thresholding small coefficients and re-fitting only the
surviving terms.  This converges in a small number of iterations and scales
well to the multi-equation setting (separate STLSQ per state variable).
The paper provides convergence arguments and demonstrates robustness to
moderate measurement noise via the implicit denoising effect of
least-squares on noisy data.

\textbf{Demonstration on benchmark dynamical systems.}
The paper demonstrates SINDy on three progressively complex systems:
(1)~the three-state Lorenz chaotic system with 7 nonzero terms, where
SINDy exactly recovers the governing equations; (2)~a nonlinear structural
oscillator with a cubic stiffness term; and (3)~a chemical reaction network
(Lorenz-type oscillator in concentration space).  In each case, SINDy
recovers the correct sparsity pattern and coefficients with high accuracy,
even at several percent measurement noise.  The cylinder wake mean-field
model studied in the companion paper by Noack et al.\ (2003) is a canonical
fluid-mechanics extension of the same framework.

\textbf{Generality and connection to existing methods.}
Brunton et al.\ argue that SINDy subsumes several existing data-driven
methods as special cases: DMD is recovered when the library contains only
linear terms; NARMAX corresponds to a polynomial library without the
sparsity constraint; and the Koopman operator perspective motivates
non-polynomial library choices.  This unifying view, combined with the
open-source implementation, has made SINDy a widely used tool and spurred
a large body of follow-on work in physics-informed machine learning.

% ============================================================
\section*{3. Replication}
% ============================================================

% ------------------------------------------------------------
\subsection*{3a. Coding Implementation}
% ------------------------------------------------------------

The SINDy methodology was implemented from scratch in Python across five
purpose-built packages totalling approximately 7{,}700 lines of source code.
No existing SINDy library (e.g.\ PySINDy) was used; every component was
written to match the algorithmic description in Brunton et al.\ (2016).

\subsubsection*{Mathematical formulation}

\paragraph{Problem statement.}
Let $\mathbf{x}(t) \in \mathbb{R}^{s}$ be the state of a dynamical system
observed at $n$ time steps, collected row-wise into
$\mathbf{X} \in \mathbb{R}^{n \times s}$.  The corresponding time-derivative
matrix $\dot{\mathbf{X}} \in \mathbb{R}^{n \times s}$ is either measured
directly or estimated from $\mathbf{X}$ via a Savitzky--Golay filter.
SINDy assumes the dynamics admit a sparse representation in a fixed
function library $\boldsymbol{\Theta}(\mathbf{X}) \in \mathbb{R}^{n \times p}$:
\begin{equation}
  \dot{\mathbf{X}} = \boldsymbol{\Theta}(\mathbf{X})\,\boldsymbol{\Xi},
  \qquad \boldsymbol{\Xi} \in \mathbb{R}^{p \times s},
  \label{eq:sindy}
\end{equation}
where $\boldsymbol{\Xi}$ is a sparse coefficient matrix.  Each column
$\boldsymbol{\xi}_{j}$ of $\boldsymbol{\Xi}$ gives the active library
terms for the $j$-th state equation $\dot{x}_{j}$.

\paragraph{Feature library.}
The library $\boldsymbol{\Theta}$ is constructed by \texttt{sindy/library.py}
from polynomial monomials up to degree $p_{\max}$ with maximum interaction
order $q_{\max}$ (number of distinct state variables per term).  For a state
$\mathbf{x} = (x_{1}, \ldots, x_{s})^{\top}$ the library columns are
\begin{equation}
  \boldsymbol{\Theta}(\mathbf{X}) =
  \bigl[\,\mathbf{1} \;\big|\; \mathbf{X} \;\big|\;
  x_{i}x_{j} \;\big|\; x_{i}^{2} \;\big|\; \cdots \bigr],
  \label{eq:library}
\end{equation}
where the constant column provides an intercept and higher-degree columns
are included subject to $|\alpha| \leq p_{\max}$ and at most $q_{\max}$
distinct indices per monomial.  For the examples in this report,
$p_{\max} = 3$ and $q_{\max} = 2$ for the Lorenz and cylinder cases,
and $p_{\max} = 2$, $q_{\max} = 2$ for the oscillator.
Near-constant columns (coefficient of variation below $10^{-3}$) are
dropped before regression to avoid numerical near-singularity.

\paragraph{Sequential Thresholded Least Squares (STLSQ).}
The core sparse solver, implemented in \texttt{AdaptiveSTLSQ}
(\texttt{sindy/fit.py}), finds $\boldsymbol{\Xi}$ for a fixed sparsity
threshold $\lambda > 0$ as follows.

\begin{enumerate}
  \item \textbf{Initialise.}  Solve the ridge-regularised normal equations
    for the full library:
    \begin{equation}
      \boldsymbol{\Xi}^{(0)} =
      \bigl(\boldsymbol{\Theta}^{\top}\boldsymbol{\Theta}
            + \alpha\mathbf{I}\bigr)^{-1}
      \boldsymbol{\Theta}^{\top}\dot{\mathbf{X}},
      \label{eq:ridge-init}
    \end{equation}
    where $\alpha$ is the ridge penalty (typically $10^{-3}$).

  \item \textbf{Threshold.}  Set to zero all coefficients with
    $|\xi^{(\ell)}_{ij}| \leq \lambda$, defining an active support
    $\mathcal{S}^{(\ell)} = \{(i,j) : |\xi^{(\ell)}_{ij}| > \lambda\}$.

  \item \textbf{Re-fit.}  For each target equation $j$ independently,
    solve a new ridge problem restricted to the active columns
    $\boldsymbol{\Theta}_{\mathcal{S}^{(\ell)}_j}$:
    \begin{equation}
      \boldsymbol{\xi}^{(\ell+1)}_{j}
      \big|_{\mathcal{S}^{(\ell)}_j}
      =
      \bigl(\boldsymbol{\Theta}_{\mathcal{S}}^{\top}
             \boldsymbol{\Theta}_{\mathcal{S}}
            + \alpha\mathbf{I}\bigr)^{-1}
      \boldsymbol{\Theta}_{\mathcal{S}}^{\top}
      \dot{\mathbf{x}}_{j}.
      \label{eq:stlsq-refit}
    \end{equation}

  \item \textbf{Iterate} steps 2--3 until
    $\|\boldsymbol{\Xi}^{(\ell+1)} - \boldsymbol{\Xi}^{(\ell)}\|_{\infty}
    < 10^{-8}$ or a maximum of 20 iterations.

  \item \textbf{Debiasing.}  On the final support, replace the Ridge solution
    with ordinary least squares (plain \texttt{lstsq}) to remove the
    $\ell_{2}$-shrinkage bias introduced by $\alpha$.
\end{enumerate}

The Ridge penalty in steps 1 and 3 stabilises the solve when library columns
are nearly collinear (e.g.\ $x$, $z$, and $xz$ in the Lorenz system are
strongly correlated along the attractor) without imposing true sparsity.
Sparsity comes solely from the hard threshold in step 2.

\paragraph{Threshold sweep and Pareto analysis.}
Rather than fixing $\lambda$ in advance, STLSQ is run for each of
$N_{\lambda} = 50$ threshold values on a logarithmic grid,
\begin{equation}
  \lambda_{i} = \lambda_{\min} \cdot
  \left(\frac{\lambda_{\max}}{\lambda_{\min}}\right)^{i/(N_{\lambda}-1)},
  \quad i = 0, \ldots, N_{\lambda}-1,
  \label{eq:lambda-grid}
\end{equation}
with $\lambda_{\min} = 10^{-3}$ and $\lambda_{\max} = 1.0$ for the examples
here.  This produces a family of candidate models
$\{(\hat{k}_i,\, \widehat{\mathrm{MSE}}_i)\}_{i=1}^{N_\lambda}$, where
$\hat{k}_i = \|\boldsymbol{\Xi}_i\|_0$ is the sparsity (number of nonzero
coefficients) and
\begin{equation}
  \widehat{\mathrm{MSE}}_i
  = \frac{1}{n \cdot s}
  \bigl\|\dot{\mathbf{X}} - \boldsymbol{\Theta}\boldsymbol{\Xi}_i\bigr\|_{F}^{2}
  \label{eq:mse}
\end{equation}
is the mean squared residual.  The \emph{Pareto-optimal front}
$\mathcal{P} \subseteq \{1,\ldots,N_\lambda\}$ retains only non-dominated
models: model $i$ is on the front if no other model $j$ has both
$\hat{k}_j \leq \hat{k}_i$ and $\widehat{\mathrm{MSE}}_j \leq
\widehat{\mathrm{MSE}}_i$ with at least one strict inequality.

\paragraph{Model selection from the Pareto front.}
Three strategies are implemented in \texttt{sindy/\allowbreak pareto.py},
applied after the sweep:

\begin{itemize}
  \item \textbf{BIC minimum.}  All $N_\lambda$ candidate models (not only the
    Pareto front) are scored with
    \begin{equation}
      \mathrm{BIC}_i = n\ln(\mathrm{MSE}_{\mathrm{eff},i})
                       + \hat{k}_i\ln(n),
      \label{eq:bic}
    \end{equation}
    where $\mathrm{MSE}_{\mathrm{eff},i} = \max\!\bigl(\widehat{\mathrm{MSE}}_i,\,
    \delta\bigr)$ and the floor $\delta = \max(10^{-3}\,\hat{\sigma}^{2}_Y,\,
    10^{-12})$ prevents the log from diverging on near-perfect fits.
    $\hat{\sigma}^{2}_Y$ is the mean per-target variance of $\dot{\mathbf{X}}$.
    The model minimising~\eqref{eq:bic} is selected.

  \item \textbf{Pareto knee.}  The point of maximum curvature on the
    log-NMSE--complexity curve, computed as the point of greatest perpendicular
    distance from the line connecting the sparsest and densest Pareto points.

  \item \textbf{Pareto dial.}  A continuous parameter $d \in [0, 1]$
    interpolates between maximum sparsity ($d = 0$) and minimum MSE ($d = 1$)
    by minimising $\ln(\mathrm{NMSE}) + (1-d)\,\lambda_{\mathrm{max}}\,
    \tilde{k}$ over the front, where $\tilde{k}$ is complexity normalised
    to $[0,1]$.  Setting $d = 1$ selects the densest Pareto point.
\end{itemize}

\paragraph{Bootstrap ensemble.}
For the Lorenz ensemble run, STLSQ is repeated $B = 50$ times, each on a
bootstrap subsample of $\lfloor 0.8n \rfloor$ rows drawn with replacement.
Each draw independently runs the full threshold sweep and selects a BIC-best
model.  The final coefficient matrix is the element-wise median across all
$B$ draws:
\begin{equation}
  \hat{\boldsymbol{\Xi}}_{\mathrm{ens}}
  = \operatorname{median}_{b=1}^{B}\bigl(\boldsymbol{\Xi}^{(b)}\bigr),
  \label{eq:ensemble}
\end{equation}
and each entry's \emph{inclusion probability}
$\pi_{ij} = B^{-1}\sum_{b=1}^{B}\mathbf{1}[|\xi^{(b)}_{ij}|>0]$
measures term stability across draws.

\subsubsection*{Code structure}

\textbf{\texttt{sindy/}} is the core package.  \texttt{library.py} builds
$\boldsymbol{\Theta}$ as described above.  \texttt{fit.py} contains
\texttt{AdaptiveSTLSQ} (the STLSQ solver), the ensemble wrapper, and
pre-processing utilities including collinear-column pruning.
\texttt{pareto.py} implements the sweep, Pareto-front extraction, and all
three selection strategies.  \texttt{pipeline.py} wraps these into a single
\texttt{run\_sindy\_pipeline\_general} call.

\textbf{\texttt{idtools/}} provides validation and diagnostics.
\texttt{compare\_to\_truth} projects the symbolic true equations (via SymPy)
onto the same library $\boldsymbol{\Theta}$ to obtain the ground-truth
coefficient matrix $\boldsymbol{\Xi}_{\mathrm{true}}$, then reports
coefficient-level relative errors and the Pearson correlation
$\rho(\boldsymbol{\xi}_{\mathrm{true}}, \boldsymbol{\xi}_{\mathrm{disc}})$,
distinguishing a good data fit (high $R^{2}$) from genuine structure
recovery (low coefficient error).

\textbf{\texttt{lorenz/}}, \textbf{\texttt{damped\_oscillator/}}, and
\textbf{\texttt{cylinder\_flow/}} each contain a self-contained ODE simulator,
SymPy truth model, SINDy configuration, and forward-integration utilities.

Primary dependencies: \textsc{NumPy} (matrix algebra),
\textsc{SciPy} (\texttt{solve\_ivp}, \texttt{lstsq}),
\textsc{scikit-learn} (\texttt{Ridge}),
\textsc{SymPy} (symbolic truth models),
\textsc{Matplotlib} (figures).

% ------------------------------------------------------------
\subsection*{3b. Replication of Simple, Verification Example}
% ------------------------------------------------------------

\subsubsection*{System}

The verification example is a 2D damped linear oscillator,
\begin{equation}
  \dot{x} = \mu x + \omega y, \qquad \dot{y} = -\omega x + \mu y,
  \label{eq:oscillator}
\end{equation}
with $\mu = -0.1$ and $\omega = 2.0$.  This system admits an exact closed-form
solution from initial condition $(x_{0}, 0)$:
\begin{equation}
  x(t) = x_{0}\,e^{\mu t}\cos(\omega t), \qquad
  y(t) = -x_{0}\,e^{\mu t}\sin(\omega t).
  \label{eq:exact}
\end{equation}
The trajectory is a decaying spiral (Figure~\ref{fig:osc-phase}): amplitude
decays at rate $e^{-0.1t}$ while oscillating at $\omega = 2\,\mathrm{rad/s}$.
The analytical solution~\eqref{eq:exact} provides ground truth at every time
step, allowing the numerical integrator and the SINDy-identified model to be
verified independently.

\subsubsection*{Why this is a suitable verification case}

The true governing equations are known exactly and are structurally simple:
only four of the twelve possible coefficients in a degree-2 polynomial library
($1, x, y, x^{2}, xy, y^{2}$ for each of two equations) are nonzero.
Correct sparse identification requires the algorithm to zero out all eight
nonlinear candidates and retain exactly $k = 4$ terms with their correct
values.

\subsubsection*{Numerical setup}

A reference trajectory was integrated over $t \in [0, 30]$ with 3{,}000
samples using SciPy's \texttt{solve\_ivp}
(RK45, rtol\,$= 10^{-10}$, atol\,$= 10^{-12}$).
Comparison to the analytical solution~\eqref{eq:exact} confirmed numerical
accuracy to within $1.8 \times 10^{-10}$ in the $\ell^{\infty}$ norm.
Exact time derivatives were supplied to SINDy by evaluating the known
right-hand side, matching the Brunton-style clean-derivative validation
used throughout the paper.

\subsubsection*{SINDy configuration and results}

SINDy was run with a degree-2 polynomial library (6 features per equation),
identity state scaling, no column normalisation, and BIC-based Pareto
selection over 50 threshold values.  The BIC selection plot
(Figure~\ref{fig:osc-bic}) shows the BIC score curve~\eqref{eq:bic}
reaching a clear minimum at $k = 4$, the correct model complexity.

The discovered equations were
\begin{equation}
  \dot{x} = -0.1\,x + 2.0\,y, \qquad \dot{y} = -2.0\,x - 0.1\,y,
  \label{eq:osc-disc}
\end{equation}
matching the true coefficients exactly.

\begin{table}[H]
  \centering
  \caption{Verification metrics for the damped oscillator example.}
  \label{tab:osc}
  \begin{tabular}{lc}
    \toprule
    Metric & Value \\
    \midrule
    Fit $R^{2}$                                          & $1.000\,000$ \\
    Coefficient correlation (true vs.\ discovered)       & $1.0000$ \\
    Max relative error on nonzero true terms             & $2.1 \times 10^{-11}$ \\
    Selected model complexity $k$                        & 4 (correct) \\
    Forward integration RMSE                             & $\approx 10^{-8}$ \\
    \bottomrule
  \end{tabular}
\end{table}

\begin{figure}[H]
  \centering
  \begin{subfigure}[t]{0.44\textwidth}
    \centering
    \includegraphics[width=\textwidth]{phase_portrait}
    \caption{Phase portrait: reference data, SINDy forward integration, and
             analytical solution~\eqref{eq:exact} overlaid.}
    \label{fig:osc-phase}
  \end{subfigure}
  \hfill
  \begin{subfigure}[t]{0.54\textwidth}
    \centering
    \includegraphics[width=\textwidth]{time_series}
    \caption{Time series $x(t)$ and $y(t)$: reference (blue), SINDy (dashed
             orange), and analytical solution (dotted green).}
    \label{fig:osc-ts}
  \end{subfigure}
  \caption{Damped oscillator trajectories.  Reference, discovered, and
           analytical curves are visually indistinguishable.}
  \label{fig:osc-traj}
\end{figure}

\begin{figure}[H]
  \centering
  \includegraphics[width=0.55\textwidth]{bic_selection}
  \caption{BIC model selection for the damped oscillator.
           \textbf{Top:} NMSE vs.\ complexity Pareto frontier; grey dots are
           all 50 threshold-sweep candidates, the blue line connects
           Pareto-optimal points, and the red star marks the BIC-selected
           model.
           \textbf{Bottom:} BIC score~\eqref{eq:bic} vs.\ complexity.
           The minimum at $k = 4$ matches the true model sparsity exactly.}
  \label{fig:osc-bic}
\end{figure}

\begin{figure}[H]
  \centering
  \begin{subfigure}[t]{\textwidth}
    \includegraphics[width=\textwidth]{structure_verification_Xi_true}
    \caption{True coefficient matrix $\boldsymbol{\Xi}_{\text{true}}$ (projected
             from symbolic RHS onto the polynomial library).}
  \end{subfigure}\\[6pt]
  \begin{subfigure}[t]{\textwidth}
    \includegraphics[width=\textwidth]{structure_verification_Xi_disc}
    \caption{Discovered coefficient matrix $\boldsymbol{\Xi}_{\text{disc}}$.}
  \end{subfigure}\\[6pt]
  \begin{subfigure}[t]{\textwidth}
    \includegraphics[width=\textwidth]{structure_verification_Xi_err}
    \caption{Coefficient error $\boldsymbol{\Xi}_{\text{disc}} -
             \boldsymbol{\Xi}_{\text{true}}$ (physical units).
             All entries are at numerical zero ($\lesssim 10^{-11}$).}
  \end{subfigure}
  \caption{Structure verification heatmaps for the damped oscillator.
           The two matrices are visually identical; the error panel is
           uniformly blank, confirming machine-precision coefficient recovery.}
  \label{fig:osc-xi}
\end{figure}

% ------------------------------------------------------------
\subsection*{3c. Replication of Real Examples}
% ------------------------------------------------------------

Two examples from Brunton et al.\ (2016) were replicated: the Lorenz system
(main paper, Figures~2--3) and the cylinder wake mean-field model (main paper,
Figure~4).

% ---- Lorenz ----
\subsubsection*{Lorenz system}

\paragraph{Setup.}
The Lorenz equations,
\begin{equation}
  \dot{x} = \sigma(y - x), \qquad
  \dot{y} = x(\rho - z) - y, \qquad
  \dot{z} = xy - \beta z,
  \label{eq:lorenz}
\end{equation}
were integrated with the paper's parameters ($\sigma = 10$, $\rho = 28$,
$\beta = 8/3$) from initial condition $(-8, 8, 27)$ over $t \in [0, 100]$
at $\Delta t = 0.001$, yielding 100{,}001 samples on the chaotic attractor.
Exact derivatives were supplied by evaluating~\eqref{eq:lorenz} directly.
SINDy used a degree-3 polynomial library with identity scaling and no column
normalisation, matching the paper's approach.

\paragraph{Results.}
The discovered equations were
\begin{align}
  \dot{x} &= -10.0\,x + 10.0\,y, \notag \\
  \dot{y} &= -xz + 28.0\,x - y, \label{eq:lorenz-disc} \\
  \dot{z} &= xy - 2.667\,z, \notag
\end{align}
recovering all seven true nonzero coefficients with no spurious terms.

\begin{table}[H]
  \centering
  \caption{Verification metrics for the Lorenz system.}
  \label{tab:lorenz}
  \begin{tabular}{lc}
    \toprule
    Metric & Value \\
    \midrule
    Fit $R^{2}$                                          & $1.000\,000$ \\
    Coefficient correlation (true vs.\ discovered)       & $1.0000$ \\
    Max relative error on nonzero true terms             & $6 \times 10^{-13}$ \\
    Selected complexity $k$                              & 7 (correct) \\
    \bottomrule
  \end{tabular}
\end{table}

\begin{figure}[H]
  \centering
  \begin{subfigure}[t]{0.48\textwidth}
    \centering
    \includegraphics[width=\textwidth]{lorenz_3d_true_vs_sindy}
    \caption{3D phase portrait: reference trajectory (blue) and SINDy
             forward integration (orange) from the same initial condition.}
    \label{fig:lorenz-3d}
  \end{subfigure}
  \hfill
  \begin{subfigure}[t]{0.48\textwidth}
    \centering
    \includegraphics[width=\textwidth]{bic_selection}
    \caption{BIC selection: the minimum at $k = 7$ matches the true number
             of nonzero Lorenz coefficients exactly.}
    \label{fig:lorenz-bic}
  \end{subfigure}
  \caption{Lorenz system replication.}
  \label{fig:lorenz}
\end{figure}

\begin{figure}[H]
  \centering
  \begin{subfigure}[t]{\textwidth}
    \includegraphics[width=\textwidth]{structure_verification_Xi_true}
    \caption{$\boldsymbol{\Xi}_{\text{true}}$}
  \end{subfigure}\\[6pt]
  \begin{subfigure}[t]{\textwidth}
    \includegraphics[width=\textwidth]{structure_verification_Xi_disc}
    \caption{$\boldsymbol{\Xi}_{\text{disc}}$}
  \end{subfigure}\\[6pt]
  \begin{subfigure}[t]{\textwidth}
    \includegraphics[width=\textwidth]{structure_verification_Xi_err}
    \caption{Error $\boldsymbol{\Xi}_{\text{disc}} - \boldsymbol{\Xi}_{\text{true}}$
             (all entries $\lesssim 10^{-12}$).}
  \end{subfigure}
  \caption{Lorenz coefficient heatmaps.  The identified $\sigma$, $\rho$,
           $\beta$, and bilinear ($xz$, $xy$) terms match the true values
           to machine precision.}
  \label{fig:lorenz-xi}
\end{figure}

% ---- Cylinder wake ----
\subsubsection*{Cylinder wake mean-field model}

\paragraph{Setup.}
The cylinder wake is described in three reduced coordinates $(x, y, z)$
corresponding to the two leading POD amplitudes and the shift mode, governed
by the mean-field model of Noack et al.\ (2003) (Equation~8 in Brunton et al.):
\begin{align}
  \dot{x} &= \mu\,x - \omega\,y + A\,xz, \notag\\
  \dot{y} &= \omega\,x + \mu\,y + A\,yz, \label{eq:cyl}\\
  \dot{z} &= -\lambda\bigl(z - x^{2} - y^{2}\bigr). \notag
\end{align}
Parameters: $\mu = 0.1$, $\omega = 1.0$, $A = -1.0$, $\lambda = 10.0$.
A reference trajectory of 3{,}000 samples over $T \in [0, 200]$ was integrated.
The notebook additionally exercises the full POD + shift-mode recovery pipeline
(\texttt{cylinder\_flow/snapshot\_pod\_shift.py}): the 3D trajectory was
embedded in an 80-dimensional random subspace and the three modal coordinates
recovered by POD; the recovered coordinates correlated with the true
coordinates at $|r| \geq 0.98$ for all three modes.

\paragraph{BIC failure and model-selection rationale.}
BIC~\eqref{eq:bic} is computed with an MSE floor,
$\mathrm{MSE}_{\mathrm{eff}} = \max(\mathrm{MSE},\; \delta)$,
where $\delta = 10^{-3}\,\mathrm{Var}(Y) \approx 3.2\times10^{-5}$.
The floor is intended to prevent a near-zero MSE from driving BIC to
$-\infty$ on clean synthetic data.  In this example, however, it has an
unintended consequence: the $k=14$ model achieves
$\mathrm{MSE} \approx 10^{-32}$ (exact recovery), which the floor clamps to
$3.2\times10^{-5}$, the same effective MSE as the $k=2$ model
($\mathrm{MSE} = 1.1\times10^{-5}$, also below the floor).
Because BIC sees both models as equally accurate, it selects the sparser
one, $k=2$, returning
\begin{equation}
  \dot{x} = -y, \qquad \dot{y} = x, \qquad \dot{z} = 0,
  \label{eq:cyl-bic}
\end{equation}
with $R^{2} = 0.66$.  Equation~\eqref{eq:cyl-bic} describes a pure
circular rotation: it captures the dominant oscillatory mode but discards
the linear growth rate $\mu$, the amplitude-limiting nonlinearity $Axz$,
$Ayz$, and the entire shift-mode equation.  It is physically incomplete.

This is a known shortcoming of BIC (and AIC) in the SINDy context: when
the residual of the dense model falls below the noise or regularisation
floor, the criterion loses discriminating power and over-prunes
\cite{Brunton2016}.  The standard diagnostic response is to inspect the
full Pareto curve rather than trusting a single scalar criterion.  Three
reference points are informative:
\begin{itemize}
  \item \textbf{BIC minimum} ($k=2$): sparsest model judged ``good enough''
        by the information criterion.  Use as a lower bound on complexity.
  \item \textbf{Pareto knee} ($k \approx 8$--$10$ here): the point of
        maximum curvature on the NMSE-vs-complexity curve, where additional
        terms yield diminishing accuracy returns.  Often a robust default.
  \item \textbf{Densest Pareto point} ($k=14$): the model at the right end
        of the frontier, with the lowest achievable MSE.  Use when physical
        knowledge constrains the expected term count, or when the knee and
        BIC minimum are implausibly sparse.
\end{itemize}
Here, both the knee ($R^{2} \approx 1$ achieved around $k=10$) and physical
knowledge of the three-mode mean-field structure suggest $k\gg 2$.
The configuration therefore sets \texttt{pareto\_dial}\,$= 1.0$, which
selects the densest Pareto point ($k=14$), and is equivalent to choosing
the minimum-MSE model regardless of complexity.  The BIC curve computed
post-hoc (Figure~\ref{fig:cyl-bic}) illustrates the collapse of the BIC
scores at $k \geq 12$ as all models hit the MSE floor, and shows visually
why the knee or the densest point must be consulted when BIC saturates.

\paragraph{Results.}
With \texttt{pareto\_dial}\,$= 1.0$ the discovered equations were
\begin{align}
  \dot{x} &= 0.1\,x - y - xz, \notag\\
  \dot{y} &= x + 0.1\,y - yz, \label{eq:cyl-disc}\\
  \dot{z} &= 10.0\,x^{2} + 10.0\,y^{2} - 10.0\,z, \notag
\end{align}
correctly recovering all structural terms of Equation~\eqref{eq:cyl} with
14 total nonzero coefficients.

\begin{table}[H]
  \centering
  \caption{Comparison of BIC-selected ($k=2$) and dial-selected ($k=14$)
           models for the cylinder wake.}
  \label{tab:cyl}
  \begin{tabular}{lccc}
    \toprule
    & BIC minimum ($k=2$) & Densest Pareto ($k=14$) & True model \\
    \midrule
    Selected by           & \texttt{pick\_mode='bic'} & \texttt{dial=1.0} & --- \\
    $R^{2}$               & $0.663$ & $1.000$ & --- \\
    $\dot{z}$ identified? & No ($\dot{z}=0$) & Yes & Yes \\
    Coeff.\ correlation   & --- & $1.0000$ & --- \\
    Physical completeness & Rotation only & Full Eq.~\eqref{eq:cyl} & Full Eq.~\eqref{eq:cyl} \\
    \bottomrule
  \end{tabular}
\end{table}

\begin{figure}[H]
  \centering
  \begin{subfigure}[t]{0.48\textwidth}
    \centering
    \includegraphics[width=\textwidth]{cylinder_3d_true_vs_sindy}
    \caption{3D phase portrait in $(x, y, z)$ POD coordinates: reference
             (blue) and SINDy integration (orange).}
    \label{fig:cyl-3d}
  \end{subfigure}
  \hfill
  \begin{subfigure}[t]{0.48\textwidth}
    \centering
    \includegraphics[width=\textwidth]{bic_selection}
    \caption{BIC curve (post-hoc).  BIC scores collapse for $k \geq 12$
             as all models hit the MSE floor; BIC selects $k=2$
             (Eq.~\eqref{eq:cyl-bic}, $R^{2}=0.66$).
             The densest Pareto point $k=14$ is used instead.}
    \label{fig:cyl-bic}
  \end{subfigure}
  \caption{Cylinder wake replication.}
  \label{fig:cyl}
\end{figure}

\begin{figure}[H]
  \centering
  \includegraphics[width=\textwidth]{cylinder_timeseries_data_vs_brunton_vs_sindy}
  \caption{Cylinder wake time series over $T \in [0, 200]$: reference data
           (blue), direct re-integration of the true Equation~\eqref{eq:cyl}
           (green dashed), and SINDy forward integration (orange dotted).
           All three agree closely throughout the window, replicating
           Figure~4 of Brunton et al.\ (2016).}
  \label{fig:cyl-ts}
\end{figure}

\begin{figure}[H]
  \centering
  \begin{subfigure}[t]{\textwidth}
    \includegraphics[width=\textwidth]{structure_verification_Xi_true}
    \caption{$\boldsymbol{\Xi}_{\text{true}}$}
  \end{subfigure}\\[6pt]
  \begin{subfigure}[t]{\textwidth}
    \includegraphics[width=\textwidth]{structure_verification_Xi_disc}
    \caption{$\boldsymbol{\Xi}_{\text{disc}}$}
  \end{subfigure}\\[6pt]
  \begin{subfigure}[t]{\textwidth}
    \includegraphics[width=\textwidth]{structure_verification_Xi_err}
    \caption{Error $\boldsymbol{\Xi}_{\text{disc}} - \boldsymbol{\Xi}_{\text{true}}$.}
  \end{subfigure}
  \caption{Cylinder wake coefficient heatmaps.  The $\mu$, $\omega$, $A$
           (bilinear $xz$, $yz$), and $\lambda$ terms of
           Equation~\eqref{eq:cyl} are all recovered correctly.}
  \label{fig:cyl-xi}
\end{figure}

\section{Github Repository}
\url{https://github.com/ksommerkohrt/SINDY_Brunton}

\begin{thebibliography}{51}
\bibitem{Brunton2016}
S.~L.\ Brunton, J.~L.\ Proctor, and J.~N.\ Kutz,
``Discovering governing equations from data by sparse identification of
nonlinear dynamical systems,''
\textit{Proc.\ Natl.\ Acad.\ Sci.}, vol.~113, no.~15, pp.~3932--3937, 2016.

\bibitem{Noack2003}
B.~R.\ Noack, K.~Afanasiev, M.~Morzy\'{n}ski, G.~Tadmor, and F.~Thiele,
``A hierarchy of low-dimensional models for the transient and post-transient
cylinder wake,''
\textit{J.\ Fluid Mech.}, vol.~497, pp.~335--363, 2003.

\bibitem{Zhang2025}
H. Zhang et al., ``Physics-Informed Machine Learning in Design and Manufacturing: Status and Challenges,'' \textit{Journal of Computing and Information Science in Engineering}, vol. 25, no. 4, 2025.

\bibitem{Kutz2023}
J. N. Kutz and S. L. Brunton, \textit{Sparse Methods for Data-Driven Modeling and Control}. Cambridge: Cambridge University Press, 2023.

\bibitem{Takeishi2026}
N. Takeishi and Y. Kawahara, ``Sparse Identification of Nonlinear Dynamics With Library Optimization Mechanism: Recursive Long-Term Prediction Perspective,'' \textit{IEEE Transactions on Cybernetics}, vol. 56, no. 5, pp. 2475--2488, 2026.

\bibitem{Morales2026}
K. Morales et al., ``Sparse identification of nonlinear dynamics and Koopman operators with Shallow Recurrent Decoder Networks,'' \textit{Proceedings of the National Academy of Sciences}, vol. 123, no. 19, 2026.

\bibitem{Vaddireddy2026}
H. Vaddireddy and A. Rasheed, ``SINDy-KANs: Sparse identification of non-linear dynamics through Kolmogorov-Arnold networks,'' \textit{arXiv preprint arXiv:2603.18548}, 2026.

\bibitem{Zhang2019}
H. Zhang and H. Schaeffer, ``On the convergence of the self-adaptive hard thresholding algorithm for sparse identification of nonlinear dynamics,'' \textit{SIAM Journal on Applied Dynamical Systems}, vol. 18, no. 2, pp. 715--744, 2019.

\end{thebibliography}

\end{document}
